\documentclass{rvcepaper}
% unit-wise question bank in the RVCE paper class
\newcommand{\qbhead}[2]{%
  \noindent
  \begin{center}{\fontsize{12}{14}\selectfont DEPARTMENT OF}\\[3pt]
  {\bfseries\fontsize{15.5}{18}\selectfont INDUSTRIAL ENGINEERING \& MANAGEMENT}\end{center}\vspace{-3pt}
  \noindent\begingroup\renewcommand{\arraystretch}{1.9}\bfseries\fontsize{11.5}{13}\selectfont
  \begin{tabular}{|p{6.2cm}|C{5.0cm}|p{5.5cm}|}\hline
   Semester \hspace{0.1cm}: \hspace{0.2cm}III & UG & Course Code \hspace{0.1cm}: \hspace{0.2cm}IM233AI\\\hline
   \multicolumn{2}{|p{11.2cm}|}{Course Title: WORK SYSTEMS DESIGN} & Question Bank\\\hline
   \multicolumn{3}{|c|}{UNIT #1 \hspace{0.3cm}--\hspace{0.3cm} #2}\\\hline
  \end{tabular}\endgroup\par\vspace{6pt}}
\newenvironment{qbtable}{%
  \begin{longtable}{|C{1.05cm}|p{10.7cm}|C{0.85cm}|C{0.95cm}|C{0.85cm}|C{0.85cm}|}
  \hline \textbf{Sl.\newline No} & \multicolumn{1}{c|}{\textbf{Questions}} & \textbf{M} & \textbf{DL} & \textbf{BT} & \textbf{CO}\\ \hline
  \endfirsthead
  \hline \textbf{Sl.\newline No} & \multicolumn{1}{c|}{\textbf{Questions}} & \textbf{M} & \textbf{DL} & \textbf{BT} & \textbf{CO}\\ \hline
  \endhead
}{\end{longtable}}
\newcommand{\qbpart}[1]{\multicolumn{6}{|c|}{\bfseries #1}\\ \hline}
\newcommand{\qbtopic}[1]{\multicolumn{6}{|l|}{\itshape #1}\\ \hline}
\newcommand{\qbrow}[6]{#1 & #2 & #3 & #4 & #5 & #6\\ \hline}
\newcommand{\qbkey}{\par\vspace{6pt}\centerline{\small M -- Marks, DL -- Difficulty Level (L: Low, M: Medium, H: High), BT -- Bloom's Taxonomy Level, CO -- Course Outcome}\par}
\begin{document}

% ---------------- Unit I
\begin{iempaper}
\paperinfo{dept={INDUSTRIAL ENGINEERING \& MANAGEMENT},title={WORK SYSTEMS DESIGN},code={IM233AI}}
\qbhead{I}{Introduction to Work Systems and Productivity}
\begin{qbtable}
\qbpart{Part -- A \hspace{0.5cm} Short-Answer Questions (2 Marks each)}
\qbrow{1}{Define the Pyramidal Structure of Work.}{2}{L}{1}{1}
\qbrow{2}{Define a Work System as a physical entity.}{2}{L}{1}{1}
\qbrow{3}{How is Work Systems defined as a field of professional practice?}{2}{L}{1}{1}
\qbrow{4}{Distinguish between a Task and a Work Element}{2}{L}{2}{1}
\qbrow{5}{What is the primary objective of Work Measurement?}{2}{L}{1}{1}
\qbrow{6}{Provide the basic formula for Productivity and define its terms.}{2}{L}{1}{1}
\qbrow{7}{List two common ways to improve productivity in a manufacturing system.}{2}{L}{2}{1}
\qbrow{8}{What is a Manual Work system?}{2}{L}{1}{1}
\qbrow{9}{What are Allowances in the context of standard time? Give one example.}{2}{L}{2}{1}
\qbrow{10}{What is meant by a Worker-Machine System?}{2}{L}{1}{1}
\qbrow{11}{Give two examples of an operation carried out using a Worker-Machine System.}{2}{L}{2}{1}
\qbrow{12}{Define an Automated Work System.}{2}{M}{1}{1}
\qbrow{13}{What is the main difference between a Worker-Machine System and an Automated System regarding the role of the worker?}{2}{M}{2}{1}
\qbrow{14}{What is the formula for Normal Time in time study? Define the terms.}{2}{M}{1}{1}
\qbrow{15}{In a worker-machine system, define the Machine Cycle Time.}{2}{L}{1}{1}
\qbrow{16}{For a single worker operating a single machine, what determines the Cycle Time?}{2}{M}{2}{1}
\qbrow{17}{State the relationship between Work Methods and Work Measurement}{2}{M}{2}{1}
\qbrow{18}{What is a Therblig? Who proposed this concept?}{2}{M}{2}{1}
\qbrow{19}{In the context of the pyramidal structure, what are Basic Motion Elements and why are they important?}{2}{M}{2}{1}
\qbrow{20}{What is Observed Time and mention how it is obtained in a time study.}{2}{M}{2}{1}
\qbpart{Part -- B \hspace{0.5cm} Descriptive Questions (8 Marks each)}
\qbtopic{Introduction to Work Systems and Productivity}
\qbrow{1}{Explain the pyramidal structure of work as described by Groover. How does it aid in work system design?}{08}{M}{3}{1}
\qbrow{2}{Discuss the importance of time in work systems. How is time used as a basis for productivity measurement?}{08}{M}{3}{1}
\qbrow{3}{Define physical work systems. Describe their components with examples.}{08}{M}{2}{1}
\qbrow{4}{How is work system a field of professional practice in industrial engineering? Illustrate with real-world applications.}{08}{M}{3}{1}
\qbrow{5}{Differentiate between manual work systems and worker-machine systems. Provide examples and implications for productivity.}{08}{M}{3}{1}
\qbrow{6}{Describe the role of standard time in work measurement. How is it established and used in planning?}{08}{M}{2}{3}
\qbrow{7}{Explain the concept of productivity improvement through work system redesign. Provide two techniques with examples.}{08}{M}{3}{1}
\qbrow{8}{Analyze the three categories of excess non-productive activities generally carried out in organizations.}{08}{M}{4}{1}
\qbrow{9}{Distinguish between Production and Productivity. Also justify the rationale for considering labor productivity as a measure of organizational efficiency and effectiveness.}{08}{M}{3}{1}
\qbrow{10}{A work group of 5 workers in a certain month produced 500 units of output working 8 hr/day for 22 days in the month.\newline (a) What productivity measures could be used for this situation, and what are the values of their respective productivity ratios?\newline (b) Suppose that in the next month, the same work group produced 600 units but there were only 20 workdays in the month. Using the same productivity measures as before, determine the productivity index using the prior month as a base.}{08}{M}{3}{1}
\qbrow{11}{A work group of 10 workers in a certain month produced 7200 units of output working 8 hr/day for 22 days in the month. Determine the labor productivity ratio using:\newline (a) Units of output per worker-hour and\newline (b) Units of output per worker-month.\newline (c) Suppose that in the next month, the same work group produced 6800 units but, there were only 20 workdays in the month. For each productivity measure in (a) and (b), determine the productivity index for the next month using the prior month as a base.}{08}{M}{3}{1}
\qbrow{12}{A work group of 20 workers in a certain month produced 8600 units of output working 8 hr/day for 21 days.\newline (a) What is the labor productivity ratio for this month?\newline (b) In the next month, the same work group produced 8000 units but there were 22 workdays in the month and the size of the work group was reduced to 14 workers. What is the labor productivity ratio for this second month?\newline (c) What is the productivity index using the first month as a base?}{08}{M}{3}{1}
\qbtopic{Manual Work and Worker-Machine Systems}
\qbrow{13}{Define a manual work system. Explain its key components and give two examples.}{08}{L}{2}{1}
\qbrow{14}{Compare manual work systems and worker-machine systems in terms of productivity, ergonomics, and cost.}{08}{M}{4}{1}
\qbrow{15}{Describe the elements of a worker-machine system. How does it differ from a fully automated system?}{08}{M}{2}{1}
\qbrow{16}{Describe the migration strategy from manual to automated work systems. What are the key considerations during this transition?}{08}{M}{2}{1}
\qbrow{17}{If the normal time is 1.30 min for a repetitive task that produces one work unit per cycle, and the company uses a PFD allowance factor of 12\%, determine:\newline (a) The standard time for the task and\newline (b) How many work units are produced in an 8-hour shift at standard performance?}{8}{L}{2}{3}
\qbrow{18}{The normal time for a repetitive task that produces two work units per cycle is 3.0 min. The plant uses a PFD allowance factor of 15\%. Determine\newline (a) The standard time per piece and\newline (b) How many work units are produced in an 8-hour shift at standard performance?}{8}{L}{2}{3}
\qbrow{19}{The normal time to perform a certain manual work cycle is 3.47 min. In addition, an irregular work element whose normal time is 3.70 min must be performed every 10 cycles. One work unit is produced each cycle. The PFD allowance factor is 14\%. Determine\newline (a) The standard time per piece and\newline (b) How many work units are produced during an 8-hour shift at 100\% performance, and the worker works exactly 7.018 hr, which corresponds to the 14\% allowance factor.\newline (c) If the worker's pace is 120\% and he works 7.2 hours during the regular shift, how many units are produced?}{8}{M}{3}{3}
\qbrow{20}{The normal time to perform a repetitive manual assembly task is 4.25 min. In addition, an irregular work element whose normal time is 1.75 min must be performed every 8 cycles. Two work units are produced each cycle. The PFD allowance factor is 16\%. Determine\newline (a) The standard time per piece and\newline (b) How many work units are produced in an 8-hour shift at standard performance?\newline (c) Determine the anticipated amount of time worked and the amount of time lost per 8-hour shift that corresponds to the PFD allowance factor of 16\%.}{8}{M}{3}{3}
\qbrow{21}{The standard time for a manual material-handling work cycle is 2.58 min per piece. The PFD allowance factor used to set the standard was 13\%. During a particular 8-hour shift of interest, it is known that the worker lost a total of 53 min due to personal time, rest breaks, and delays. On that same day, the worker completed 214 work units. Determine\newline (a) The number of standard hours accomplished,\newline (b) Worker efficiency, and\newline (c) The worker's performance level expressed as a percentage.}{8}{M}{3}{3}
\qbrow{22}{The normal time of the work cycle in a worker--machine system is 5.39 min. The operator-controlled portion of the cycle is 0.84 min. One work unit is produced each cycle. The machine cycle time is constant. (a) Using a PFD allowance factor of 16\% and a machine allowance factor of 30\%, determine the standard time for the work cycle. (b) If a worker assigned to this task completes 85 units during an 8-hour shift, what is the worker's efficiency? (c) If it is known that a total of 42 min was lost during the 8-hour clock time due to personal needs and delays, what was the worker's performance on the portion of the cycle he controlled?}{8}{M}{3}{3}
\qbrow{23}{A worker is responsible for loading and unloading a production machine. The load/unload elements in the repetitive work cycle have a normal time of only 24 sec, and the machine cycle time is 2.83 min. One part is produced each cycle. Every sixth cycle, the operator must replace the tote pans of parts, which takes 2.40 min (normal time). For setting the standard time, the PFD allowance factor is 15\% and the machine allowance factor is 15\%. Determine the standard time under the following alternative assumptions: (a) the irregular element is performed as an external element and (b) the irregular element is performed as an internal element. (c) Determine the corresponding standard daily production quantities (8-hour shift) for each of these time standards.}{8}{M}{4}{3}
\qbrow{24}{The work cycle in a worker--machine system consists of (1) external manual work elements with a total normal time of 0.42 min, (2) a machine cycle with a machine time of 1.12 min, and (3) internal manual elements with a total normal time of 1.04 min. (a) Determine the standard time for the cycle, using a PFD allowance factor of 15\% and a machine allowance factor of 30\%. (b) How many work units are produced daily (8-hour shift) at standard performance?}{8}{M}{3}{1}
\qbrow{25}{The normal time for a work cycle in a worker--machine system is 6.27 min. For setting the standard time, the PFD allowance factor is 12\% and the machine allowance factor is 25\%. The work cycle includes manual elements totaling a normal time of 5.92 min, all but 0.65min of which are performed as internal elements. Determine (a) the standard time for the cycle and (b) the daily output at standard performance. (c) During an 8- hour shift, the worker lost 39 min due to personal time, rest breaks, and delays, and she produced 72 pieces. What was the worker's pace on the operator-controlled portion of the shift?}{8}{M}{4}{3}
\end{qbtable}
\qbkey
\end{iempaper}

% ---------------- Unit II
\begin{iempaper}
\paperinfo{dept={INDUSTRIAL ENGINEERING \& MANAGEMENT},title={WORK SYSTEMS DESIGN},code={IM233AI}}
\qbhead{II}{Methods Engineering and Operations Analysis}
\begin{qbtable}
\qbpart{Part -- A \hspace{0.5cm} Short-Answer Questions (2 Marks each)}
\qbrow{1}{Define Methods Engineering.}{2}{L}{1}{2}
\qbrow{2}{What is the main objective of Methods Engineering?}{2}{L}{1}{2}
\qbrow{3}{Distinguish between Method Study and Work Measurement.}{2}{M}{2}{2}
\qbrow{4}{What is a Fishbone Diagram (Cause-and-Effect Diagram) and when is it used?}{2}{M}{2}{2}
\qbrow{5}{What are the five standard symbols used in a Flow Process Chart?}{2}{L}{1}{2}
\qbrow{6}{Distinguish between a Flow Process Chart and an Operation Process Chart (OPC).}{2}{M}{2}{2}
\qbrow{7}{What information does a Flow Diagram convey?}{2}{L}{1}{2}
\qbrow{8}{What is the primary purpose of a Multi-Activity Chart?}{2}{L}{2}{2}
\qbrow{9}{What is a Block Diagram used for in methods engineering?}{2}{L}{2}{2}
\qbrow{10}{How does a Process Map differ from a Flow Process Chart?}{2}{L}{2}{2}
\qbrow{11}{State two main categories under which the Principles of Motion Economy are grouped.}{2}{L}{1}{2}
\qbrow{12}{According to the Principles of Motion Economy, why should hands begin and complete their motions simultaneously?}{2}{M}{2}{2}
\qbrow{13}{Define the Therbligs- Reach and Move.}{2}{L}{1}{2}
\qbpart{Part -- B \hspace{0.5cm} Descriptive Questions (8 Marks each)}
\qbtopic{Introduction to Methods Engineering \& Operations Analysis}
\qbrow{1}{Explain the role of Methods Engineering in improving productivity and work design.}{8}{M}{2}{2}
\qbrow{2}{List and explain the steps in the methods engineering procedure.}{8}{M}{2}{2}
\qbrow{3}{Explain the relationship between methods engineering and automation. How does methods analysis support automation decisions?}{8}{M}{4}{2}
\qbrow{4}{Explain how methods engineering contributes to lean manufacturing and continuous improvement.}{8}{M}{4}{2}
\qbrow{5}{How does the ``USA Principle'' of automation apply to Methods Engineering?}{8}{M}{3}{2}
\qbtopic{Charting and Diagramming Techniques}
\qbrow{6}{Explain how flow process charts are used in methods engineering. Provide an example of how they help improve a manual task.}{8}{M}{3}{2}
\qbrow{7}{Explain the difference between operation process charts and flow process charts. Provide one example of each.}{8}{M}{4}{2}
\qbrow{8}{Describe how block diagrams are used in methods engineering. What are their advantages in system visualization?}{8}{M}{2}{2}
\qbrow{9}{Discuss the role of traditional IE charting techniques in identifying non-value-added activities. Provide two examples.}{8}{M}{4}{2}
\qbrow{10}{Explain the construction and use of an ``Operation Process Chart.''}{8}{M}{2}{2}
\qbrow{11}{Describe the ``Flow Diagram'' (Spatial) and its relationship to the Flow Process Chart.}{8}{M}{3}{2}
\qbrow{12}{Explain the use of ``Two-Handed Process Chart'' (Right-hand/Left-hand chart) in the application of Principles of Motion Economy.}{8}{M}{3}{2}
\qbrow{13}{What is a Man-Machine Chart (Multiple Activity Chart), and how is it used to calculate efficiency?}{8}{M}{2}{2}
\qbrow{14}{Discuss ``Check Sheets'' and ``Pareto Analysis'' as data collection tools in Methods Engineering.}{8}{M}{2}{2}
\qbrow{15}{Review any four types of basic Data Collection and Analysis techniques in Methods Engineering.}{8}{M}{2}{2}
\qbrow{16}{Explain and illustrate the three stages of Automation Migration Strategy.}{8}{M}{3}{2}
\qbrow{17}{Explain any four techniques of Methods Engineering.}{8}{M}{2}{2}
\qbtopic{Principles of Motion Economy}
\qbrow{18}{Discuss how principles of motion economy are applied in methods engineering. Provide two examples.}{8}{H}{4}{2}
\qbrow{19}{Define the principles of motion economy. How do they contribute to efficient work design?}{8}{M}{2}{2}
\qbrow{20}{Explain any six principles of motion economy related to the use of the human body and explain each briefly.}{8}{M}{2}{2}
\qbrow{21}{Explain how workplace layout affects motion economy. Suggest three layout improvements for a manual assembly station.}{8}{M}{3}{2}
\qbrow{22}{Identify and explain three principles of motion economy related to the arrangement of the workplace.}{8}{M}{2}{2}
\qbrow{23}{Discuss the principles of motion economy specifically related to the use of the two hands during a manual task.}{8}{M}{3}{2}
\qbrow{24}{``The hands should be relieved of work elements that can be performed by other body parts.'' Explain this principle with respect to jigs, fixtures, and foot pedals.}{8}{M}{3}{2}
\qbrow{25}{Describe the concept of ``Normal'' and ``Maximum'' Reach Zones in workstation design.}{8}{M}{2}{2}
\end{qbtable}
\qbkey
\end{iempaper}

% ---------------- Unit III
\begin{iempaper}
\paperinfo{dept={INDUSTRIAL ENGINEERING \& MANAGEMENT},title={WORK SYSTEMS DESIGN},code={IM233AI}}
\qbhead{III}{Work Measurement and Time Study}
\begin{qbtable}
\qbpart{Part -- A \hspace{0.5cm} Short-Answer Questions (2 Marks each)}
\qbrow{1}{Define the Standard Time for a task.}{2}{L}{1}{3}
\qbrow{2}{State the fundamental formula used to calculate Standard Time from Normal Time.}{2}{L}{1}{3}
\qbrow{3}{Define Normal Performance (or Normal Pace) in the context of time standards.}{2}{L}{1}{3}
\qbrow{4}{What is Performance Rating (PR) in time study?}{2}{L}{1}{3}
\qbrow{5}{List two key prerequisites for determining a valid time standard using the stop-watch method.}{2}{L}{1}{3}
\qbrow{6}{Why must the work method be standardized before setting a time standard?}{2}{M}{2}{3}
\qbrow{7}{Name the three major types of Allowances (PFD) included in Standard Time.}{2}{L}{1}{1}
\qbrow{8}{Define Personal Needs Allowance and give one example of its use.}{2}{M}{2}{1}
\qbrow{9}{Distinguish between Basic Fatigue Allowance and Variable Fatigue Allowance.}{2}{L}{2}{3}
\qbrow{10}{How is the total Allowance usually expressed when calculating Standard Time?}{2}{L}{1}{1}
\qbrow{11}{List the four main methods for determining time standards, ranked by relative accuracy (Highest to Lowest).}{2}{M}{2}{3}
\qbrow{12}{List four main methods for determining time standards, ranked by relative application speed (Fast -Slow).}{2}{M}{2}{3}
\qbrow{13}{Why is the PMTS method generally considered the most accurate?}{2}{M}{2}{3}
\qbrow{14}{What is the purpose of selecting a ``Qualified Worker'' at the start of the time study procedure?}{2}{M}{2}{3}
\qbrow{15}{Define the term Cycle Time in the context of time study.}{2}{L}{1}{1}
\qbrow{16}{What does the acronym MTM stand for, and what is the unit of time used in MTM-1?}{2}{L}{1}{3}
\qbrow{17}{What is the fundamental concept behind the simplified PMTS known as MOST ?}{2}{L}{1}{3}
\qbrow{18}{List the three main sequence models used in Basic MOST.}{2}{L}{1}{3}
\qbrow{19}{Write the standard formula (sequence) for the General Move sequence in Basic MOST.}{2}{L}{1}{3}
\qbrow{20}{What is Mini MOST and for what type of work is it primarily used?}{2}{M}{2}{3}
\qbrow{21}{What is Maxi MOST and when is it typically applied?}{2}{M}{2}{3}
\qbrow{22}{What is the purpose of MOST for Windows (or electronic PMTS systems)?}{2}{L}{2}{3}
\qbpart{Part -- B \hspace{0.5cm} Descriptive Questions (8 Marks each)}
\qbtopic{Work Measurement and Time Study}
\qbrow{1}{Define a ``Standard Time'' according to Groover and explain the four specific conditions that must be met for a time standard to be valid.}{8}{L}{2}{3}
\qbrow{2}{Rank the five primary work measurement methods by accuracy and justify the high ranking of Predetermined Motion Time Systems (PMTS).}{8}{M}{2}{3}
\qbrow{3}{Rank the work measurement methods by ``Relative Application Speed'' (Time to set the standard) and explain why Historical Data is the fastest.}{8}{M}{2}{3}
\qbtopic{Allowances and Time Study}
\qbrow{4}{Distinguish between Personal, Fatigue, and Delay (PFD) allowances. Why are they necessary in determining the final standard time?}{8}{M}{2}{1}
\qbrow{5}{Explain the difference between ``Constant Fatigue Allowances'' and ``Variable Fatigue Allowances.'' Provide examples of factors that influence variable fatigue.}{8}{M}{3}{1}
\qbrow{6}{What are ``Contingency Allowances''? How do they differ from ``Policy Allowances''?}{8}{M}{2}{1}
\qbrow{7}{A worker's observed time is 2.0 min, and their performance rating is 110\%. The company uses a PFD allowance of 15\%. Calculate the Normal Time and the Standard Time.}{8}{M}{3}{1}
\qbrow{8}{Explain the ``Internal'' vs. ``External'' work elements in a worker-machine system. How do allowances apply differently to these?}{8}{M}{2}{1}
\qbrow{9}{Outline the step-by-step procedure for determining the Standard Time of a manual task starting from the initial observation.}{8}{M}{2}{3}
\qbrow{10}{Why is it necessary to break a task into ``Work Elements'' before timing it? Provide ANY six rules for defining these elements.}{8}{M}{2}{3}
\qbrow{11}{An operator takes an average of 1.20 minutes to complete a cycle. The analyst rates the worker at 115\%. If one element in the cycle was a machine-controlled process taking 0.40 minutes, calculate the Normal Time for the manual portion and the total cycle.}{8}{M}{3}{3}
\qbrow{12}{Discuss the procedure for determining standard time in ``Machine-Paced'' tasks. How do allowances differ for the machine portion vs. the manual portion?}{8}{M}{2}{3}
\qbrow{13}{The snap back timing method was used to obtain average times for work elements in one work cycle. The times are given below. Element d is a machine-controlled element and the time is constant. Elements a, b, c, e and f are operator controlled and were performance rated at 80\%; however elements e and f are performed during the machine controlled element d. The machine allowance is zero and the PFD allowance is 14\%. Determine the normal time and standard time for the cycle.\par \qtabc{|l|c|c|c|c|c|c|}{\hline Element & a & b & c & d & e & f\\\hline Observed Time (min) & 0.24 & 0.30 & 0.17 & 0.76 & 0.26 & 0.14\\\hline }}{8}{M}{3}{3}
\qbrow{14}{The continuous timing method in direct time study was used to obtain the elemental times for a worker machine task as indicated. Element c is a machine controlled and the time is constant. Elements a, b, d, e and f are operator controlled and external to machine cycle and they are performance rated at 80\%. If the machine allowance is 25\% and worker PFD allowance is 15\%, determine the normal time and standard time for the work cycle.\par \qtabc{|l|c|c|c|c|c|c|}{\hline Element & a & b & c & d & e & f\\\hline Observed Time (min) & 0.18 & 0.30 & 0.88 & 1.12 & 1.55 & 1.80\\\hline }}{8}{M}{3}{3}
\qbrow{15}{The continuous timing method in direct time study was used to obtain the elemental times for a worker machine task as indicated. Element c is a machine controlled and the time is constant. Elements a, b, d, e and f are operator controlled and external to machine cycle and they are performance rated at 90\%. If the machine allowance is 12\% and worker PFD allowance is 8\%, determine the normal time and standard time for the work cycle.\par \qtabc{|l|c|c|c|c|c|c|}{\hline Element & a & b & c & d & e & f\\\hline Observed Time (min) & 0.24 & 0.38 & 0.88 & 1.04 & 1.65 & 1.90\\\hline }}{8}{M}{3}{3}
\qbrow{16}{The continuous timing method in direct time study was used to obtain the elemental times for a worker machine task as indicated. Element c is a machine controlled and the time is constant. Elements a, b, d, e and f are operator controlled and external to machine cycle and they are performance rated at 80\%. If the machine allowance is 25\% and worker PFD allowance is 15\%, determine the normal time and standard time for the work cycle.\par \qtabc{|l|c|c|c|c|c|c|}{\hline Element & a & b & c & d & e & f\\\hline Observed Time (min) & 0.12 & 0.28 & 0.78 & 1.04 & 1.45 & 1.80\\\hline }}{8}{M}{3}{3}
\qbrow{17}{The continuous timing method in direct time study was used to obtain the elemental times for a worker machine task as indicated. Element c is a machine controlled and the time is constant. Elements a, b, d, e and f are operator controlled and external to machine cycle and they are performance rated at 90\%. If the machine allowance is 35\% and worker PFD allowance is 25\%, determine the normal time and standard time for the work cycle.\par \qtabc{|l|c|c|c|c|c|c|}{\hline Element & a & b & c & d & e & f\\\hline Observed Time (min) & 0.12 & 0.28 & 0.78 & 1.04 & 1.45 & 1.80\\\hline }}{8}{M}{3}{3}
\qbrow{18}{The continuous timing method in direct time study was used to obtain the elemental times for a worker machine task as indicated. Element c is a machine controlled and the time is constant. Elements a, b, d, e and f are operator controlled and external to machine cycle and they are performance rated at 90\%. If the machine allowance is 35\% and worker PFD allowance is 25\%, determine the normal time and standard time for the work cycle.\par \qtabc{|l|c|c|c|c|c|c|}{\hline Element & a & b & c & d & e & f\\\hline Observed Time (min) & 0.18 & 0.30 & 0.88 & 1.12 & 1.55 & 1.80\\\hline }}{8}{M}{3}{3}
\qbrow{19}{The flyback timing method was used to obtain average times for work elements in one work cycle. The times are given below. Element d is a machine-controlled element and the time is constant. Elements a, b, c, e and f are operator controlled and were performance rated at 90\%; however elements e and f are performed during the machine controlled element d. The machine allowance is zero and the PFD allowance is 14\%. Determine the normal time and standard time for the cycle.\par \qtabc{|l|c|c|c|c|c|c|}{\hline Element & a & b & c & d & e & f\\\hline Observed Time (min) & 0.18 & 0.24 & 0.21 & 0.74 & 0.26 & 0.24\\\hline }}{8}{M}{3}{3}
\qbrow{20}{The snap back timing method was used to obtain average times for work elements in one work cycle. The times are given below. Element d is a machine-controlled element and the time is constant. Elements a, b, c, e and f are operator controlled and were performance rated at 90\%; however elements e and f are performed during the machine controlled element d. The machine allowance is zero and the PFD allowance is 14\%. Determine the normal time and standard time for the cycle.\par \qtabc{|l|c|c|c|c|c|c|}{\hline Element & a & b & c & d & e & f\\\hline Observed Time (min) & 0.24 & 0.30 & 0.19 & 0.96 & 0.26 & 0.24\\\hline }}{8}{M}{3}{3}
\qbtopic{Predetermined Motion Time Systems (PMTS)}
\qbrow{21}{Discuss the advantages of PMTS over direct stopwatch time study as highlighted by Groover.}{8}{L}{2}{3}
\qbrow{22}{Distinguish between first-level (MTM-1) and higher-level (MTM-2, MTM-3) PMT systems.}{8}{M}{2}{3}
\qbrow{23}{What is MOST (Maynard Operation Sequence Technique)? Explain its focus on ``Sequence Models'' rather than individual motions.}{8}{M}{2}{3}
\qbrow{24}{Break down the ``General Move'' sequence model (ABGABPA). Define each parameter.}{8}{M}{2}{3}
\qbrow{25}{Compare ``General Move'' vs. ``Controlled Move'' in Basic MOST. Give an example for each.}{8}{M}{3}{3}
\qbrow{26}{Distinguish between Mini-MOST, Basic MOST, and Maxi-MOST based on cycle length and repetitiveness.}{8}{M}{2}{3}
\qbrow{27}{What is ``Admin MOST''? Provide three examples of tasks where it would be applied.}{8}{M}{3}{3}
\qbrow{28}{Discuss ``MOST for Windows'' (Computerized MOST). What are the advantages of using software over manual data cards?}{8}{M}{2}{3}
\qbrow{29}{Outline the systematic procedure for applying a Predetermined Motion Time System (PMTS) to establish a time standard for a manual task.}{8}{L}{2}{3}
\qbrow{30}{Describe the ``Tool Use'' sequence model. List four common tool-use actions it covers.}{8}{M}{2}{3}
\end{qbtable}
\qbkey
\end{iempaper}

% ---------------- Unit IV
\begin{iempaper}
\paperinfo{dept={INDUSTRIAL ENGINEERING \& MANAGEMENT},title={WORK SYSTEMS DESIGN},code={IM233AI}}
\qbhead{IV}{Lean Production, Waste Elimination and Worker Involvement}
\begin{qbtable}
\qbpart{Part -- A \hspace{0.5cm} Short-Answer Questions (2 Marks each)}
\qbrow{1}{What is the Japanese term for Waste in Lean Production? Define it.}{2}{L}{2}{2}
\qbrow{2}{Name the two most detrimental forms of waste identified in Lean production.}{2}{L}{1}{2}
\qbrow{3}{Why is Overproduction considered the worst form of waste?}{2}{M}{2}{2}
\qbrow{4}{Define the Lean waste known as Defects (or Production of defective parts).}{2}{L}{1}{2}
\qbrow{5}{Define the Lean waste known as Waiting.}{2}{L}{1}{2}
\qbrow{6}{List two ``other forms of waste'' (non-Muda wastes) often discussed in Lean.}{2}{L}{1}{2}
\qbrow{7}{How does the Lean concept of Elimination of Waste contribute to cost reduction?}{2}{M}{2}{2}
\qbrow{8}{Define the core principle of Just-in-Time (JIT) Production.}{2}{L}{1}{2}
\qbrow{9}{What is a Pull System of Production Control?}{2}{L}{1}{2}
\qbrow{10}{What is the role of Kanban in JIT production?}{2}{M}{1}{2}
\qbrow{11}{Distinguish between an Internal Setup and an External Setup.}{2}{M}{2}{2}
\qbrow{12}{What is the Japanese term for Autonomation? Define its core concept.}{2}{L}{1}{2}
\qbrow{13}{State the two essential elements of Jidoka.}{2}{L}{1}{2}
\qbrow{14}{What does Total Productive Maintenance (TPM) focus on?}{2}{L}{1}{2}
\qbrow{15}{Define the Japanese term for Continuous Improvement.}{2}{L}{1}{2}
\qbrow{16}{List the five Ss in the 5S methodology (in English).}{2}{L}{1}{2}
\qbrow{17}{Define Takt Time and explain its significance in Lean.}{2}{M}{2}{2}
\qbpart{Part -- B \hspace{0.5cm} Descriptive Questions (8 Marks each)}
\qbtopic{Lean and Waste}
\qbrow{1}{Explain the ``Seven Wastes'' (Muda) in a production system with specific focus on defective parts and overproduction.}{8}{M}{2}{2}
\qbrow{2}{Define ``Muda'' and explain the waste of Overproduction. Why is it considered the most serious waste?}{8}{M}{4}{2}
\qbrow{3}{Compare and contrast ``Waste of Motion'' and ``Waste of Transportation.'' Provide one industrial example for each and evaluate how excessive inventory contributes to both.}{8}{H}{4}{2}
\qbrow{4}{Describe the elements of a typical structure of Lean Manufacturing System. Support your answer with an illustration.}{8}{M}{2}{2}
\qbrow{5}{Define Over-processing as a form of waste. Provide an example from a machining context.}{8}{M}{3}{2}
\qbrow{6}{Explain the waste of Waiting (Idle time). What are typical causes in a manufacturing plant?}{8}{M}{2}{2}
\qbtopic{Just in Time Production}
\qbrow{7}{Compare and contrast Push vs. Pull systems of production control. Why is Pull preferred in Lean?}{8}{H}{4}{2}
\qbrow{8}{Explain the SMED (Single Minute Exchange of Die) methodology for setup time reduction.}{8}{M}{3}{2}
\qbrow{9}{Why is ``Setup Time Reduction'' a prerequisite for implementing a Pull system?}{8}{H}{4}{2}
\qbrow{10}{Describe the function of a ``Kanban'' card in a dual-card system with an example.}{8}{M}{3}{2}
\qbrow{11}{Explain the Pull System of production control using the Kanban mechanism.}{8}{M}{2}{2}
\qbrow{12}{Analyze the Internal vs. External Setup activities in the context of SMED.}{8}{M}{4}{2}
\qbtopic{Autonomation}
\qbrow{13}{What is Total Productive Maintenance (TPM)? Explain the goal of ``Zero Breakdowns.''}{8}{M}{2}{2}
\qbrow{14}{Explain the concept of Poka-Yoke (Error Proofing) with three industrial manufacturing examples.}{8}{H}{3}{2}
\qbrow{15}{What is Jidoka? Explain how ``Stopping the Process'' actually improves long-term efficiency.}{8}{H}{4}{2}
\qbrow{16}{Explain the ``Andon'' system and its role in the Jidoka philosophy.}{8}{M}{3}{2}
\qbrow{17}{What is Total Productive Maintenance (TPM)? List its primary goals regarding machine availability.}{8}{M}{2}{2}
\qbtopic{Worker Involvement}
\qbrow{18}{Define ``Standardized Work'' and explain its three essential components according to Groover.}{8}{M}{2}{2}
\qbrow{19}{Discuss the role of ``Takt Time'' in aligning production rate with customer demand.}{8}{M}{3}{1}
\qbrow{20}{Explain the 5S System of visual management and its impact on workplace productivity.}{8}{M}{3}{1}
\qbrow{21}{Analyze the organizational requirements for a successful Continuous Improvement culture?}{8}{H}{4}{2}
\end{qbtable}
\qbkey
\end{iempaper}

% ---------------- Unit V
\begin{iempaper}
\paperinfo{dept={INDUSTRIAL ENGINEERING \& MANAGEMENT},title={WORK SYSTEMS DESIGN},code={IM233AI}}
\qbhead{V}{Ergonomics, Human Factors and Occupational Safety}
\begin{qbtable}
\qbpart{Part -- A \hspace{0.5cm} Short-Answer Questions (2 Marks each)}
\qbrow{1}{Define Ergonomics. What is its primary goal?}{2}{L}{2}{4}
\qbrow{2}{State the relationship between Ergonomics and Human Factors.}{2}{M}{2}{4}
\qbrow{3}{Name the three main thematic/ topic areas or sub-disciplines of Ergonomics.}{2}{L}{1}{4}
\qbrow{4}{What is the focus of Physical Ergonomics?}{2}{L}{2}{4}
\qbrow{5}{What is the focus of Cognitive Ergonomics? Give an example.}{2}{M}{2}{4}
\qbrow{6}{List the three basic components that form a Human-Machine System.}{2}{L}{1}{4}
\qbrow{7}{Describe the ``Interactions'' that occur within a Human-Machine System.}{2}{M}{2}{4}
\qbrow{8}{Name two potential negative effects of excessive Noise exposure in the workplace.}{2}{M}{1}{4}
\qbrow{9}{Define an Industrial Accident in the context of Occupational Safety.}{2}{L}{1}{4}
\qbrow{10}{Define Occupational Disorders. Give one example.}{2}{M}{2}{4}
\qbrow{11}{What is the term for occupational disorders related to musculoskeletal strain? Give an example.}{2}{M}{1}{4}
\qbrow{12}{What is Anthropometry and how is it used in workplace design?}{2}{M}{2}{4}
\qbrow{13}{Why is the concept of ``fit the task to the man'' central to ergonomics?}{2}{M}{2}{4}
\qbrow{14}{Distinguish between static (structural) anthropometry and dynamic (functional) anthropometry.}{2}{M}{2}{4}
\qbpart{Part -- B \hspace{0.5cm} Descriptive Questions (8 Marks each)}
\qbtopic{Introduction to Ergonomics \& Human Factors}
\qbrow{1}{What is meant by the term Ergonomics and discuss its primary objectives.}{8}{L}{2}{4}
\qbrow{2}{Describe the components of a Human-Machine System.}{8}{M}{2}{4}
\qbrow{3}{Discuss the ``Human as a System Component'' including strengths and limitations.}{8}{M}{4}{4}
\qbrow{4}{Explain the three major domains (themes) of Ergonomics. Discuss how each domain contributes to the optimization of human well-being and overall system performance.}{8}{M}{4}{4}
\qbrow{5}{Explain the role of Anthropometry in work system design.}{8}{H}{4}{4}
\qbrow{6}{Describe the Human Information Processing Model.}{8}{M}{2}{4}
\qbrow{7}{Analyze the impact of Illumination on work performance.}{8}{M}{4}{4}
\qbrow{8}{Discuss Noise and its control in the workplace.}{8}{M}{4}{4}
\qbrow{9}{Discuss ``Environmental Stressors'' and their cumulative effect on productivity.}{8}{M}{4}{4}
\qbrow{10}{Discuss the ``Fitting the Task to the Man'' (FTTM) philosophy. Why is this approach superior to ``Fitting the Man to the Task'' (FMTT)?}{8}{M}{3}{4}
\qbtopic{Occupational Safety and Health}
\qbrow{11}{How can ``Poka-Yoke'' (Error Proofing) be used as an ergonomic strategy to prevent industrial accidents? Give two examples.}{8}{M}{3}{4}
\qbrow{12}{Differentiate between ``Unsafe Acts'' and ``Unsafe Conditions.'' Which is generally more prevalent in industrial accidents, and why?}{8}{M}{3}{4}
\qbrow{13}{What are Work-related Musculoskeletal Disorders (WMSDs)? List the primary risk factors.}{8}{M}{2}{4}
\qbrow{14}{Describe the ``Reach Envelope'' and its importance in workstation layout.}{8}{M}{3}{4}
\qbrow{15}{Define ``Industrial Accident'' and differentiate between ``Accident'' and ``Injury''.}{8}{M}{2}{4}
\qbrow{16}{Analyze the role of ``Unsafe Acts'' as a primary cause of accidents in workplaces.}{8}{M}{4}{4}
\qbrow{17}{Discuss the impact of ``Human Factors'' on accident rates.}{8}{M}{3}{4}
\qbrow{18}{What is meant by ``Occupational Disease'' and how it differs from an ``Industrial Injury.''}{8}{M}{2}{4}
\qbrow{19}{Develop an ``Ergonomic Checklist'' for an office worker to self-assess their risk for developing Repetitive Strain Injuries (RSI).}{8}{H}{5}{4}
\qbrow{20}{What are ``Near Misses'' and why are they critical for industrial safety management?}{8}{M}{2}{4}
\qbrow{21}{Describe ``Unsafe Conditions'' in the context of workplaces and provide four workplace examples.}{8}{M}{3}{4}
\end{qbtable}
\qbkey
\end{iempaper}

\end{document}
